\label{sec:eval}

\noindent\textbf{本节回答}：本文的方法在什么范围内可信？优点是“无冲突”的\emph{充要}编码与与求解器分离的独立校验，
使每个结果都可被逐位复核；局限是最优性只在撤销层与问题 3 被证明，调整层只有值与下界，
若干题意读法本文只能并列给出而无法判定。

\subsection{模型的优点}

\begin{enumerate}[leftmargin=1.8em,itemsep=0.25em]
  \item \textbf{统一的单元格语言}。四个问题——检测、消解、增容、放宽间隔——被统一到“时频单元集合是否相交”
        这一个判据上（命题~\ref{prop:cell}、\ref{prop:exact}）。这既使模型精确（单元约束是无冲突的\emph{充要}编码，
        不是近似或松弛），又使校验独立（同一判据的另一种算法实现即构成校验器）。
  \item \textbf{可证明的部分被明确证明}。问题 2 撤销层的全局最优性（定理~\ref{thm:q2opt}）与问题 3 的最优性
        （定理~\ref{thm:q3}）都给出了闭合的对偶间隙，后者更由三个独立上界同时确认。
        推论~\ref{cor:nozero} 直接回答了“能否一个都不撤销”这一自然追问。
  \item \textbf{未证明的部分被诚实标注}。调整各级只报告“值 / 证明下界 / 状态”（表~\ref{tab:q2solver}、\ref{tab:q4solver}），
        不作最优性声明。现任解上界割（命题~\ref{prop:cut}）在数学上保证报告值不劣于任何已知可行解，
        从而消除了“交叉验证击败主解”这类自相矛盾。
  \item \textbf{主观选择被量化}。目标函数的形式化是本文最主观的一步，四方案对照（表~\ref{tab:q2schemes}）
        把“优先级条款值多少”变成了可以看见的数字：坚持不撤销 A、B 类计划的代价，
        是撤销总数从 $\valQtwoSchemeTCancelTotal$ 上升到 $\valQtwoSchemePCancelTotal$。
  \item \textbf{题意歧义被正面处理}。问题 3 的两种读法给出相差约五倍的答案，本文不回避，
        而是以模板一致性为准选定主解释，并为另一解释给出可核查的区间式~\eqref{eq:bracket}。
  \item \textbf{结果可复现，且复现的边界被写明}。论文报告的全部数值产生于同一次云端运行，
        每个数字由阶段产物自动注入论文，每个阶段记录代码版本、参数与输入输出的 SHA-256（附录~\ref{app:repro}）；
        支撑材料另附不依赖云端账号的本地驱动脚本。唯一需要外部输入的是热启动现任解，
        本文报告了它的来源、本轮是否改进了它，以及不加载它时的冷启动结果（第~\ref{sec:q2alg} 节）。
\end{enumerate}

\subsection{模型的局限}
\label{sec:limits}

\begin{enumerate}[leftmargin=1.8em,itemsep=0.25em]
  \item \textbf{调整层的最优性未闭合}。问题 2 的调整三级与幅度级、问题 4 的撤销 C 级均未证明最优
        （\valQtwoAllOptimal{}、\valQfourAllOptimal{}）。其间隙较宽，主要原因是线性松弛在装填型约束上很弱
        （撤销级的 LP 下界为 $0$）。改进方向是加入更强的有效不等式（如基于冲突图团的覆盖不等式）或列生成。
  \item \textbf{问题 3 的答案依赖于具体的问题 2 方案}。定理~\ref{thm:q3} 的 $\valQthreeNewPlans$ 是
        “在本文这一个无冲突方案之上”的最大值。本文在 $\valQthreeLayoutCases$ 个无冲突方案（表~\ref{tab:q2schemes} 的四种读法）
        上各解一次装填，可装填数介于 $\valQthreeLayoutMin$ 与 $\valQthreeLayoutMax$ 之间（相差 $\valQthreeLayoutSpreadPercent\%$），
        而对应的自由单元数只相差 $\valQthreeLayoutFreeSpreadPercent\%$——由引理~\ref{lem:tiling}，
        可装填数量主要取决于既有布局与相位划分的\emph{对齐程度}，而非自由单元的\emph{总数}。
        若要回答“跨所有无冲突方案的最大值”，需把问题 2 与问题 3 合并为一个联合优化，规模将显著增大。
  \item \textbf{问题 4 的“不增加时频资源”是本文自加的解释}。题目只在问题 3 写明该要求，问题 4 未重复；
        本文按同一尺度加入硬约束 $e_i\le\valQfourHorizonCap$（第~\ref{sec:q4horizon} 节），
        并同时报告去掉该约束的变体（撤销 $\valQfourFreeCancelTotal$、最晚结束时刻 $\valQfourFreeHorizon$）。
        若命题人本意是不设限，应以变体为准；两种读法的结果都已给出，本文不替评阅者做这一判断。
  \item \textbf{离散化与动作集合是题目给定的}。$\Delta f$、$\Delta t$ 的粒度、“只改一个参数”的限制、
        平移幅度上限均为外生。连续化模型、多参数联合调整、功率或方向性维度都未建模。
  \item \textbf{物理层机制未建模}（假设~\ref{as:indep}）。真实电磁兼容需要保护带、空间隔离、
        互调与杂散抑制等；本文的“冲突”是纯几何的。好在模型对此是\emph{可扩展}的：
        只需把 $U(P)$ 沿频率（或时间）方向膨胀一个保护宽度，全部命题与算法原样成立，只是冲突对数会上升。
  \item \textbf{大规模的可证明最优性未验证}。合成实例在固定时限下只达到可行解（\valBenchScaleAllOptimal{}）。
\end{enumerate}

\subsection{模型的推广}

\paragraph{（一）带保护带的冲突判据。}把频段方向的占用由 $[f,f+w)$ 膨胀为 $[f-\rho,\,f+w+\rho)$
（$\rho$ 为保护频段数），时间方向同理膨胀为收发转换时间，命题~\ref{prop:cell}--\ref{prop:exact} 与全部算法不变。
这使模型可直接用于真实的频谱指配。

\paragraph{（二）在线与滚动消解。}实际用频计划是陆续提报的。把已批准的计划视为固定占用、
新提报的计划视为待放置对象，问题即退化为问题 3 的装填模型；若允许有限次回溯调整，
则是问题 2 模型在小邻域上的重解。由于检测是亚秒级的（表~\ref{tab:benchscaling}），
在线判定完全可以实时完成。

\paragraph{（三）鲁棒性与不确定性。}若计划的起始时刻带有抖动 $\pm\epsilon$，
可把每次使用的时间区间膨胀 $\epsilon$ 后重解，得到“对抖动免疫”的方案；
方案的代价随 $\epsilon$ 的增长曲线正是第~\ref{sec:sens} 节灵敏度分析的变体。

\paragraph{（四）同类问题。}本文的“每个对象恰选一个位置、位置之间不得共享资源单元、
按优先级词典序最优”的结构，同样适用于会议室与航班时刻分配、光网络波长指配、
数据中心作业的时间--机架放置等问题；单元格 AtMostOne 编码与现任解保持的词典序搜索可原样迁移。

\subsection{给用频管理部门的建议}

基于本文的定量结果，提出三条建议：
\textbf{其一}，适度放宽频段调整权限比增加带宽更有效——把最大频段平移幅度由 $10\Delta f$ 提高到 $15\Delta f$，
即可把必须撤销的计划数从 $\valSensBaseCancelTotal$ 降为 $\valSensMinCancelTotal$（第~\ref{sec:sens} 节）；
\textbf{其二}，允许 C 类装备微调用频间隔是低成本的松弛——在不增加时频资源（$e_i\le\valQfourHorizonCap$）的前提下，
它把撤销数从 $\valQtwoCancelTotal$ 降到 $\valQfourCancelTotal$（表~\ref{tab:q4compare}），
代价是多调整 $\valQfourAdjustDelta$ 个计划；
\textbf{其三}，若提报时能引导 C 类装备按引理~\ref{lem:tiling} 的相位错开（首次时刻相差 $d_C$ 的整数倍），
同一频带的利用率可接近 $100\%$，从源头上减少冲突。
